% Generated draft: requires mathematical and source review.
\documentclass[11pt]{article}
\usepackage[margin=25mm]{geometry}
\usepackage{fontspec}
\setmainfont{Latin Modern Roman}
\usepackage{amsmath,amssymb,mathtools,unicode-math}
\setmathfont{Latin Modern Math}
\usepackage{enumitem,needspace,xurl,hyperref}
\setlength{\parindent}{0pt}
\setlength{\parskip}{5pt}
\setlength{\emergencystretch}{3em}
\newcommand{\N}{\mathbb{N}}
\newcommand{\Z}{\mathbb{Z}}
\newcommand{\Q}{\mathbb{Q}}
\newcommand{\R}{\mathbb{R}}
\newcommand{\C}{\mathbb{C}}
\newcommand{\F}{\mathbb{F}}
\newcommand{\A}{\mathbb{A}}
\newcommand{\PP}{\mathbb{P}}
\newcommand{\E}{\mathbb{E}}
\newcommand{\eps}{\varepsilon}
\newcommand{\dd}{\,\mathrm{d}}
\begin{document}
\section*{UnsolvedMath ID: 30004355}
\textbf{Source problem code:} catalogue entry
\subsection{Definitions and mathematical statement}
For an undirected graph \(G=(V,E)\) whose edges have unit length, the path distance between vertices \(u,v\) in the same connected component is
\[
 d_G(u,v)=\min\{k\in\Z_{\geq 0}:\text{there is a path of }k\text{ edges from }u\text{ to }v\}.
\]
This is one possible geometric structure on a graph, included only to explain terminology; the source does not specify this interpretation or impose these graph hypotheses. It also does not specify a class of differentiable manifolds or the geometric structures on them.

The supplied research question is, verbatim:
\begin{quote}
What can be said about the global or large-scale geometry of graphs and its relationship with differential geometry on manifolds?
\end{quote}
The source does not define ``global or large-scale geometry,'' select graph or manifold classes, specify a comparison relation, or state a conclusion with quantifiers. Consequently, this record specifies a research direction rather than a precise mathematical proposition. A quantified conjecture cannot be extracted without additional hypotheses and an intended target; the illustrative distance definition above does not supply those missing choices.
\subsection{Short English statement}
The question asks how global geometric properties of graphs relate to differential geometry on manifolds. The graph classes, geometric structures, meaning of large scale, and desired relationship are unspecified, so the supplied evidence does not determine a precise conjecture.
\subsection{Sources}
\begin{itemize}
\item Supplied local UnsolvedMath registry record in \texttt{data/unsolvedmath/problems.json}, permanent ID \(30004355\), title ``Large-Scale Geometry of Graphs and Manifolds,'' source problem code \texttt{catalogue entry}; collection: Oberwolfach Reports Open Problems. The record lists proposed year \(2020\) and no proposer.
\item Supplied report citation: \emph{Combinatorics (2021)}, Oberwolfach Reports, DOI 10.4171/owr/2020/1, pp. 6--89. Supplied source URL: \url{https://doi.org/10.4171/owr/2020/1}. The background describes Oberwolfach Reports as documenting mathematical workshops and their open-problem discussions. The report contents and the location of this question within the cited page range have not been independently verified.
\item Supplied literature-source label: ``Linial, large-scale geometry of graphs,'' \url{https://doi.org/10.4171/owr/2020/1}. This label is retained as supplied; it does not establish authorship or proposer attribution for this question.
\item Literature assessment dated exactly \texttt{2026-08-22}: supplied dated curation, not independently verified. The record labels the question ``open'' and assesses it as a broad research question about large-scale graph geometry rather than a precise proposition with a known resolution. It further states that substantial theories exist but that the extracted question remains open in this form. The background describes this as dated literature triage derived from the supplied OWR curation and requests independent verification of time-sensitive status and cited claims. These are supplied assessments, not verified mathematical conclusions.
\end{itemize}
\end{document}
